fMRI data quality is severely compromised by factors such as motion, physiological effects, and thermal noise from the scanner \citep{liu2016}. Typically, fMRI data are collected with a gradient‑recalled‑echo EPI sequence at a single echo time (TE) chosen to maximise BOLD sensitivity for a given region of interest. However, T2* varies across brain regions and subjects. Multi‑echo (ME) fMRI overcomes this limitation by acquiring and combining signals at multiple TEs after a single excitation pulse, thereby optimising BOLD temporal characteristics and sensitivity \citep{kundu2012differentiating}. Furthermore, ME‑fMRI enables the separation of TE‑dependent (BOLD) from non‑TE‑dependent (non‑BOLD) signals using independent component analysis (ME‑ICA) \citep{kundu2012differentiating, dupre2021te}, which has been shown to improve test‑retest reliability with shorter scan times compared to single‑echo acquisitions \citep{lynchRapidPrecisionFunctional2020}. Even after these steps, thermal noise remains and can constrain sensitivity and specificity, especially at high spatial resolutions where thermal noise dominates.

Among thermal denoising algorithms for single‑echo fMRI, NORDIC has emerged as an effective method that operates on the complex signal, whereas MPPCA operates on the magnitude alone. However, an integrated approach combining the advantages of both TE‑dependent processing and thermal noise reduction has not yet been formulated for ME‑fMRI, despite promising alternatives such as NORDIC for VASO MRI \citep{knudsen2025} and TMPPCA for multi‑TE diffusion MRI \citep{olesenTensorDenoisingMultidimensional2023}. In this thesis, we systematically evaluate four thermal denoising strategies (MPPCA, NORDIC, HYDRA, and TMPPCA) on ME‑fMRI data across resting‑state and task‑based paradigms, using three datasets: two low‑resolution multi‑subject cohorts (2.4 mm and 2 mm isotropic) and one high‑resolution single‑subject dataset (1.5 mm) where thermal noise is dominant. We assess temporal SNR (tSNR), preservation of $S_0$ and $R_2^*$ estimates, physiological plausibility of $\Delta R_2^*$ fluctuations, functional connectivity reliability via split‑half resampling, and task‑related sensitivity using a face localiser and cross‑validated GLM prediction.

Across all analyses, TMPPCA consistently yields the highest tSNR improvements, the closest preservation of $S_0$ and $R_2^*$ to undenoised data, and the largest increase in physiologically plausible $\Delta R_2^*$ estimates. It also produces the strongest gains in functional connectivity reliability and task‑based sensitivity (higher t‑statistics in the fusiform face area and better cross‑run prediction of evoked time series). NORDIC  show moderate improvements. Notably, this is the first study to examine thermal denoising specifically in the multi‑echo fMRI domain. Our findings establish TMPPCA as a promising  for reducing thermal noise in ME‑fMRI, offering improved data quality without compromising the signal estimates fundamental to BOLD detection.